\paragraph{One benchmark, one package.}
A package does not grow a second, unrelated task family. When the question changes, that is a new package, so a result names exactly what was measured.

\paragraph{Deterministic graders first.}
Each task returns a score from 0 to 1 and a reason string. The engine can call a model judge. It does not do so unless the benchmark asks. None of the packages in this release ask. A judged score, if a later release adds one, has to be reported as a separate figure beside the deterministic score.

\paragraph{Every source carries the same provenance fields.}
Every source has a name, a kind, a URL, an SPDX licence, the steps to obtain it, a retrieval time, a contamination note, and a refresh rule. The kinds are synthetic, public-dataset, human-written, live-feed, and commissioned. The preferred order is a committed seeded generator, then a permissive public dataset, then text we wrote. We do not copy tasks out of a benchmark whose licence forbids it. Table~\ref{tab:suite} is the measured mix. Most rows are synthetic or human-written. The only live feed is the OpenRouter price snapshot used by \texttt{cost-per-task}, and that row carries licence \texttt{NOASSERTION} because the endpoint does not state an SPDX licence.

\paragraph{A hidden split is a held-out draw from a known family.}
The public split is the one in the npm package. Generators take a seed and can write a hidden split that stays out of git and out of the tarball. Publishing the generator reveals the task families, so a hidden split cannot promise novelty; the provenance notes say so, and serious comparisons should use one.

\paragraph{Episodes are traces.}
A seeded environment exposes tools, accepts calls, and can be reset. The grader replays the tool log. It does not trust a final state the model attached, and it does not trust tool results the model wrote into the log. Single-answer tasks stay first-class: \suiteSingle{} of the \suitePackages{} packages are single-answer and the other \suiteEpisode{} are episodes.

\paragraph{Cost stays missing when the provider omits it.}
The report stores token and dollar totals when the adapter sends them. The mock provider reports zero. We do not invent a price for a run that had none.

\paragraph{Write the grader's failure modes down first.}
Several packages keep a list of ways the grader can be wrong, written before the implementation, and an end-to-end script that tries the cheap attacks: empty output, random text, an always-yes answer, a repeated tool call. That does not prove the grader correct. It is a check that the obvious ways to game it do not score like a real solution. Appendix~\ref{app:suite} names one failure mode per package and the guard that watches it.
